% RELATED WORK — drafted 2026-09-08 from docs/language-resources.md and
% docs/licensing.md, both of which read licences and data-availability
% statements at source rather than from abstracts.
%
% NOTHING IN THIS FILE IS CITED FROM AN ABSTRACT. Where the repository does not
% record a bibliographic detail, that detail is NOT written here -- see the
% TODO block at the end for what still needs completing from source.

\section{Related work}
\label{sec:related}

\subsection{Why we make no comparison with a triage instrument}
\label{sec:related:instruments}

We compare against no established triage instrument, and the reason is not
licensing. It is that every instrument we examined is defined on a different kind
of input from the one this system receives, which makes a comparison not merely
unavailable but ill-posed.

The instrument our own specification named as the basis for its three categories
is the WHO Emergency Triage Assessment and Treatment (ETAT) manual for
participants (WHO, 2005, ISBN 92\,4\,154687\,5). Reading it settled the question
against us, and the details are worth stating because they are specific:

\begin{itemize}
\item \textbf{It is defined on examination signs.} ``Triage is the process of
  rapidly \emph{examining} all sick children when they first arrive in
  hospital''; ``Triage of patients involves \emph{looking for signs} of serious
  illness or injury'' (Module One, pp.~3--4). Each emergency sign is elicited by
  hand or eye: taking the child's hand to judge warmth, pressing the nail bed and
  timing capillary refill, pinching the abdominal skin, comparing the child's
  palms with the assessor's own (Modules Three to Five, pp.~26, 35--36, 44).
\item \textbf{It rules explicitly against history where history and observation
  could conflict.} On convulsions: ``This assessment depends on your observation
  of the child and not on the history from the parent \dots\ The child must be
  seen to have a convulsion during the triage process or while waiting in the
  outpatient department'' (p.~7). Where history is used --- whether the child has
  diarrhoea, a history of choking, a referral note --- it is asked of the
  caretaker in person, as an adjunct to an examination that still takes place.
\item \textbf{It does not address remote or written assessment.} Triage happens
  ``as soon as a sick child arrives in the hospital, well before any
  administrative procedure such as registration'' (p.~5). Across all 83 pages of
  the extracted text, the words \emph{telephone}, \emph{phone}, \emph{radio},
  \emph{remote}, \emph{SMS} and \emph{mobile} do not occur once.
\item \textbf{Its population is children.} ``a process of rapid triage for all
  children presenting to hospital'' (p.~1); adults appear only as equipment sizes
  (pp.~21, 46, 65). No upper age limit is stated anywhere in the document.
\item \textbf{Its categories are its own}: EMERGENCY, PRIORITY and NON-URGENT
  (p.~4; Chart~2, pp.~67--68), not the critical/urgent/routine triple our
  specification uses, and not the Emergency Severity Index levels our
  specification also cites in the same glossary entry.
\end{itemize}

The Emergency Severity Index and the Manchester Triage System are bedside
instruments of the same kind: a trained assessor observes a present patient. Our
system receives text written by a patient or a relative who has not been
examined, often before any contact with a health worker. An instrument defined on
observed signs cannot be a reference standard for a system that never observes
anything, and a number comparing the two would be measuring agreement between
incommensurable procedures. Licensing is a real constraint on reproducing their
descriptors --- the ESI grant covers clinical use of the algorithm but not
redistribution of its text, and the MTS is proprietary --- but it is the second
reason, not the first.

\paragraph{The absence is the finding.} We could not locate any validated
instrument for assigning urgency from a patient's own written description without
examination, in any language. If one exists, it is the reference standard this
work needs and we did not find it; if one does not, then the construct this
system would measure has not been established by anyone, and that is the more
important statement. We record it as an open question addressed to clinical
readers rather than as a gap we have closed, and we treat every urgency label in
our corpus as an unvalidated construct until a clinician rules otherwise.

\subsection{Clinical guidance consulted, and what we reproduce from it}

Where published guidance exists for a presentation, we consulted it while writing
concepts: the WHO IMCI Chart Booklet (2014, ISBN 978\,92\,4\,150682\,3), the
WHO--ICRC \emph{Basic Emergency Care} (2018, CC BY-NC-SA 3.0 IGO), and WHO
\emph{Managing Complications in Pregnancy and Childbirth} (2017, CC BY-NC-SA 3.0
IGO). We reproduce no text from any of them. The IMCI booklet is all rights
reserved, and the two share-alike documents would propagate their licence term
into a released corpus. A clinical fact may be cited; the wording in which a
document expresses it is not ours to copy.

We give no count of how many concepts carry such an anchor. Our own audit found
the concept total recorded with several different values across the repository,
and the anchor count with it, so any figure we quoted here would be one of
several and not a measurement. Establishing that count against the source
documents is outstanding work, and the honest position is to say so rather than
to pick one.

\subsection{Kinyarwanda language resources}

Kinyarwanda NLP has usable encoders and retrievers and very little else, and
almost nothing in a patient's voice. The resources we assessed fall into three
groups.

\paragraph{News-domain corpora.} KINNEWS and KIRNEWS provide 3{,}449 Kinyarwanda
news articles collected from a single news site. MasakhaNER 1.0 and 2.0 provide
named-entity-annotated Kinyarwanda and Swahili text, also news. Both are useful
for pretraining and neither contains clinical language or patient speech. We use
neither: MasakhaNER is non-commercial share-alike, and the KINNEWS licence is not
stated in its repository at all, which we treat as unusable until confirmed
rather than as permissive by default.

\paragraph{Patient-facing health text.} We located one collection of genuine
patient-facing Kinyarwanda health material, and it covers immunisation and
nothing else, so it maps to a single domain of our taxonomy. A Swahili health
corpus with an open licence exists and, measured directly rather than assumed,
contains institutional text --- ministerial press releases, parliamentary
addresses --- with zero lines combining a first-person marker and a symptom
word. Register, not volume, is the binding constraint on this task.

\paragraph{Community health worker questions.} The most useful Kinyarwanda
clinical resource we found is a set of real community health worker questions
with Kinyarwanda originals and clinician answers, released as a 524-question
subset under CC BY 4.0 alongside a study of large language models for frontline
healthcare support in low-resource settings. The full collection of 5{,}609
questions is not released; the data availability statement names an access route
rather than a download. We read the statement at source, and this is one of the
cases where doing so mattered: the headline figure in the abstract is not the
figure that was released.

That resource is close to ours in setting and different in kind. It records what
a health worker asks a clinician; we need what a patient says to a health worker.
It is a genuine sample of real language, which our corpus is not, and its
existence is the strongest argument that the authored approach we took is a
stopgap rather than a destination.

\subsection{Where this work sits}

We are not aware of a prior patient-voice medical triage dataset for
Kinyarwanda, and the audit above did not locate one. That is a statement about
our search, not about the field. What we can say more firmly is narrower and
more useful: the resources that exist are news-domain or institutional, the one
clinical Kinyarwanda corpus we found is in a health worker's voice rather than a
patient's, and neither the concept taxonomy nor the licensing position we
adopted could be inherited from any of them.

% ---------------------------------------------------------------------------
% TODO BEFORE SUBMISSION -- BIBLIOGRAPHY ENTRIES NOT YET VERIFIED AT SOURCE.
%
% The claims above are all substantiated by docs/language-resources.md and
% docs/licensing.md. The BIBLIOGRAPHIC DETAILS below are not all recorded
% there, and are deliberately not invented in the prose. Complete each from
% the source document itself, not from a search result or an abstract:
%
%   - CHW study. Recorded: "Large language models for frontline healthcare
%     support in low-resource settings", Nature Health 2026,
%     nature.com/articles/s44360-025-00038-1, open access via PMC12880909,
%     article CC BY 4.0. Released subset: figshare 10.6084/m9.figshare.29213147,
%     CC BY 4.0, 524 questions. NEEDED: full author list, volume/pages.
%     Recorded names are Samuel Rutunda (first author on the figshare dataset),
%     Emery Hezagira and Eric Remera (RBC co-authors), and Bilal Akhter Mateen
%     (corresponding author on the protocol paper) -- these are contact-route
%     notes, NOT a verified author list, and must not be used as one.
%   - KINNEWS / KIRNEWS. NEEDED: authors, year, venue. Recorded: 3,449
%     Kinyarwanda news articles from igihe.com; licence not stated in the repo.
%   - MasakhaNER 1.0 and 2.0. NEEDED: authors, year, venue for both. Recorded
%     licence: CC-BY-4.0-NC, with monolingual source data under per-site terms.
%   - Masakhane. Named in docs/generative-assistant-assessment.md only as a
%     source of small news-domain sets. If it is to be cited as a research
%     collective rather than as a dataset, that needs its own reference and a
%     sentence that says something specific -- do not cite it decoratively.
%   - ESI and MTS. NEEDED: citable references IF they are named in prose at all.
%     Recorded: ESI (c) Emergency Nurses Association, royalty-free for clinical
%     use of the algorithm, reproduction requires written permission; MTS
%     proprietary via ALSG. No comparison is made against either (2026-09-17).
%   - WHO ETAT manual for participants, 2005, ISBN 92 4 154687 5. Page
%     references in sec:related:instruments are printed-page numbers verified
%     in reports/TAXONOMY_SCOPE.md §2 against the PDF in docs/clinical/.
%   - WHO IMCI 2014, WHO-ICRC BEC 2018, WHO MCPC 2017. NEEDED: full titles and
%     WHO publication identifiers. IMCI's ISBN is recorded above.
%   - The immunisation-only patient-facing collection and the Swahili health
%     corpus are described but not named in the prose. Name them once the
%     entries are verified; the measurements quoted are from
%     docs/swahili-corpus-assessment.md.
% ---------------------------------------------------------------------------
